% Part: turing-machines
% Chapter: machines-computations
% Section: church-turing-thesis

\documentclass[../../../include/open-logic-section]{subfiles}

\begin{document}

\olfileid{tur}{mac}{ctt}
\olsection{تز چرچ--تورینگ}

قرار است ماشین‌های تورینگ جایگزینی دقیق برای مفهوم روش مؤثر باشند.
تورینگ بر این باور بود که هرکس هم مفهوم روش مؤثر و هم مفهوم ماشین
تورینگ را دریابد، به‌طور شهودی خواهد پذیرفت که هر کاری که با روشی
مؤثر انجام‌پذیر باشد با ماشین تورینگ نیز انجام‌پذیر است. این ادعا از
این واقعیت پشتیبانی می‌گیرد که همهٔ جایگزین‌های دقیق دیگری که برای
مفهوم روش مؤثر پیشنهاد شده‌اند، از نظر مصداقی با مفهوم ماشین تورینگ
هم‌ارز از کار درآمده‌اند؛ یعنی دقیقاً همان مجموعهٔ توابع را محاسبه
می‌کنند. این ادعا \emph{تز چرچ--تورینگ} نام دارد.

\begin{defn}[تز چرچ--تورینگ]
\emph{تز چرچ--تورینگ} می‌گوید هر چیزی که به‌وسیلهٔ روشی مؤثر
محاسبه‌پذیر باشد، تورینگ‌محاسبه‌پذیر است.
\end{defn}

تز چرچ--تورینگ به دو شیوه به کار می‌رود. کاربرد نخست آن بهانه‌ای برای
کم‌کاری است. فرض کنید توصیفی از روشی مؤثر برای محاسبهٔ چیزی، مثلاً به
صورت «شبه‌کد»، در دست داریم. آنگاه می‌توانیم به تز چرچ--تورینگ استناد
کنیم تا این ادعا را موجه سازیم که همان تابع را ماشین تورینگی محاسبه
می‌کند، حتی اگر در واقع آن ماشین را نساخته باشیم.

کاربرد دیگر تز چرچ--تورینگ از نظر فلسفی جالب‌تر است. می‌توان نشان داد
توابعی وجود دارند که ماشین‌های تورینگ نمی‌توانند آن‌ها را محاسبه
کنند. با استفاده از تز چرچ--تورینگ می‌توان از این نتیجه گرفت که آن‌ها
با هیچ روشی نمی‌توانند به‌طور مؤثر محاسبه شوند. زیرا اگر چنین روشی
وجود داشت، بنا بر تز چرچ--تورینگ برای آن ماشین تورینگی نیز وجود
می‌داشت. پس اگر ثابت کنیم هیچ ماشین تورینگی آن را محاسبه نمی‌کند،
هیچ روش مؤثری نیز وجود ندارد. به‌ویژه، از تز چرچ--تورینگ استفاده
می‌شود تا بگوییم مسئلهٔ موسوم به توقف نه‌تنها با ماشین‌های تورینگ
حل‌شدنی نیست، بلکه اساساً هیچ حل مؤثری ندارد.

\end{document}
